% Fragmento público generado desde propietarios íntegros.
% Residencia: 52.6.2.
% Fuente: ../incoming/radion_monografia_20260827/manuscrito/sections/07_incertidumbre_oscilador.tex:156-197
\paragraph{Intercambio energético entre los sectores eléctrico y magnético}

Los dos términos de \eqref{eq:HLC} se convierten en
\begin{equation}
U_E(\theta)=\hret\omega\cos^2\theta,
\qquad
U_B(\theta)=\hret\omega\sin^2\theta.
\label{eq:energy-exchange}
\end{equation}
Por tanto,
\[
U_E+U_B=\hret\omega.
\]
La polaridad eléctrica--magnética no consiste en etiquetar arbitrariamente
los focos. Consiste en dos reservas conjugadas que se intercambian durante la
órbita manteniendo constantes frecuencia y acción total.

El cociente común se prolonga:
\begin{equation}
\boxed{
e^{-\eta}
=\frac{b_S}{a_S}
=\frac{\Delta P}{\Delta Q}
=\frac{Z_\eta}{Z_*}
=\sqrt{\frac{A-C^*}{A+C^*}}.}
\label{eq:shape-chain-LC}
\end{equation}
Los productos conservados son
\[
a_Sb_S=2\hret,\qquad
\det V_\eta=\frac{\hret^2}{4},
\qquad
L_\eta C_\eta=\omega^{-2}.
\]

\begin{exactbox}
El oscilador LC es la realización física mínima de la doble lectura: el
producto conserva acción y frecuencia; el cociente conserva orientación e
impedancia. Un ciclo completo intercambia electricidad y magnetismo sin
perder el área \(\oint P\,\diff Q=\hfull\).
\end{exactbox}

